\chapter{PENUTUP}
\label{bab:penutup}

Bab ini menunjukkan penempatan simpulan dan saran dalam susunan contoh. Uraiannya merangkum fungsi template dan tidak menyatakan adanya hasil penelitian ilmiah. Ketika template dipakai untuk skripsi, ganti bagian ini dengan jawaban atas rumusan masalah berdasarkan hasil dan pembahasan yang sudah tersedia.

\section{Simpulan}

Contoh naskah menyediakan struktur enam bab, halaman awal, daftar referensi, dan lampiran. Unsur tabel, diagram vektor, persamaan, serta sitasi ditampilkan untuk memperlihatkan pola penulisan dan rujukan silang. Identitas, isi ilmiah, redaksi administratif, dan sumber yang relevan perlu disiapkan oleh pengguna sesuai penelitian serta ketentuan yang berlaku.

Tabel \ref{tab:contoh-rangkuman} memperlihatkan cara memetakan fokus dengan bentuk jawaban. Isi tabel adalah rangkuman fungsi contoh, bukan kesimpulan yang diperoleh dari pengujian. Penulis dapat memakai susunannya untuk memeriksa apakah rumusan masalah dan simpulan pada skripsi sesungguhnya mempunyai hubungan yang jelas.

\begin{table}[htbp]
  \centering
  \caption{Contoh pemetaan fokus penulisan dan bentuk rangkuman}
  \label{tab:contoh-rangkuman}
  \begin{tabular}{p{3.8cm}p{8.8cm}}
    \toprule
    Fokus penulisan & Bentuk rangkuman pada contoh \\
    \midrule
    Hubungan antarbagian & Bab dihubungkan melalui pengantar, sistematika, dan rujukan silang \\
    Unsur naskah & Tabel, gambar, persamaan, sitasi, serta lampiran mempunyai contoh yang dapat disunting \\
    Pemeriksaan dokumen & Isi dan tampilan perlu diperiksa setelah materi contoh diganti \\
    \bottomrule
  \end{tabular}
\end{table}

Dalam penelitian sesungguhnya, simpulan sebaiknya menyatakan jawaban secara langsung dan menjaga batas bukti. Pernyataan yang lebih luas daripada objek, sumber, atau kondisi yang diperiksa dapat memberi makna yang tidak didukung. Jumlah simpulan disesuaikan dengan pertanyaan serta hubungan antartemuan, sehingga bukan sekadar mengikuti jumlah poin pada contoh.

\section{Saran}

Pengguna disarankan menyiapkan garis besar isi sebelum mengganti paragraf contoh. Langkah ini membantu menyesuaikan fungsi tiap bab dengan jenis penelitian. Setelah materi diganti, rujukan silang, judul tabel, daftar referensi, dan lampiran diperiksa kembali agar tetap sesuai dengan alur serta unsur yang benar-benar tersedia.

Redaksi halaman persetujuan dan orisinalitas perlu disesuaikan dengan formulir yang berlaku. Nama dan nomor identitas harus diperiksa sebelum berkas dipakai untuk keperluan akademik. Ruang tanggal dan tanda tangan hanya diisi berdasarkan proses sebenarnya. Penggunaan template tidak menggantikan persetujuan pembimbing terhadap naskah.

Pada skripsi akhir, saran diturunkan dari keterbatasan yang telah dibahas. Saran dapat menyebut pemeriksaan tambahan, sumber yang perlu diperluas, atau keputusan metode yang perlu diuji pada kondisi lain. Penulis perlu menjelaskan alasan usulan tersebut tanpa menjadikannya sebagai janji keberhasilan penelitian lanjutan.
